\section{Anatomy of a Dock}

The clearest way to grasp the whole base model is to zoom all the way in. Look at one \emph{dock}. Look at its twenty-four \emph{sites}, the directions leading out of it. Look at the \emph{tones} sitting on those sites. Then watch one \emph{beat} act. This single picture puts the geometry of the mesh and the law of the knit together in one concrete object, and once it is clear the rest of the theory is what that object does at scale.

A short orientation first. The universe of this theory is a \textbf{vibe mesh}, a flow built from eight atoms. The relevant atoms here are the \emph{mesh} (the whole tessellation), the \emph{dock} (one cell, a here-now), the \emph{site} (one of the directions out of a dock), the \emph{tone} (the ternary value a site carries), the \emph{knit} (the law that updates the state), the \emph{wake} (the growing edge of the mesh), and the \emph{beat} (the discrete time step). This section assembles them around the dock.

\begin{principle}[The contents of a dock]
A dock holds exactly twenty-four tones, one tone per site. There is no rest slot, no twenty-fifth site, no home. Twenty-four directed ternary values are the dock's entire state, and the array of all docks' tones is the entire universe.
\end{principle}

\subsection{What one dock is}

A dock is a \textbf{here-now}, a single cell of space-time. It is the smallest place where anything happens. The space of the theory is not a container that docks float in. It is the weave of docks itself. A dock plus its neighbours is the whole local geometry, and distance is nothing more than step-counting on the graph those neighbours define.

Each dock has exactly \textbf{twenty-four neighbours}, one reached along each site. The twenty-four sites are the twenty-four $D_4$ root vectors, every coordinate arrangement of $(\pm 1, \pm 1, 0, 0)$ in four dimensions. The count is immediate. Six ways to choose which two of the four axes $x, y, z, w$ are nonzero, times four ways to choose the two signs, gives $6 \times 4 = 24$. Every site has the same length, norm squared $2$, so the sites are \textbf{single-speed}. Nothing moves faster or slower than anything else at the base.

\begin{definition}[The sites of a dock]
The twenty-four \emph{sites} of a dock are the twenty-four $D_4$ root vectors $(\pm 1, \pm 1, 0, 0)$ and all their axis rearrangements. Equivalently they are the vertices of the regular four-dimensional $24$-cell, and equivalently the unit Hurwitz quaternions. These are one set of twenty-four objects in three guises.
\end{definition}

Table~\ref{tab:the-twenty-four-sites} lists the full set, grouped by which two of the four axes are nonzero.

\begin{table}[ht]
\centering
\begin{tabular}{@{}lll@{}}
\toprule
index & nonzero axes & the four directions \\
\midrule
$0,1,2,3$ & $(x,y)$ & $(+,+,0,0)\ (+,-,0,0)\ (-,+,0,0)\ (-,-,0,0)$ \\
$4,5,6,7$ & $(x,z)$ & $(+,0,+,0)\ (+,0,-,0)\ (-,0,+,0)\ (-,0,-,0)$ \\
$8,9,10,11$ & $(x,w)$ & $(+,0,0,+)\ (+,0,0,-)\ (-,0,0,+)\ (-,0,0,-)$ \\
$12,13,14,15$ & $(y,z)$ & $(0,+,+,0)\ (0,+,-,0)\ (0,-,+,0)\ (0,-,-,0)$ \\
$16,17,18,19$ & $(y,w)$ & $(0,+,0,+)\ (0,+,0,-)\ (0,-,0,+)\ (0,-,0,-)$ \\
$20,21,22,23$ & $(z,w)$ & $(0,0,+,+)\ (0,0,+,-)\ (0,0,-,+)\ (0,0,-,-)$ \\
\bottomrule
\end{tabular}
\caption{The twenty-four sites of a dock, the $D_4$ roots, grouped by axis-pair.}
\label{tab:the-twenty-four-sites}
\end{table}

The symmetry group of the sites is $F_4$, of order $1152$, the symmetry group of the $24$-cell. It is finite. That a discrete group this large governs the local directions is what lets continuous rotation reappear only in the emergent infrared, as a coarse-grained limit, and never in the base.

That is the entire local geometry. A dock and its twenty-four neighbours. Distance is step-counting along sites on this graph.

\begin{principle}[Distance is step-counting]
There is no background space the docks sit in. The mesh is the graph of docks joined by sites, and the distance between two docks is the least number of single-site hops between them. Step-counting is what space means.
\end{principle}

\subsection{The opposite map and the lines}

For a tone to stream it needs a well-defined reverse direction. The \textbf{opposite} of a site is its negation. Site $(+,+,0,0)$ has opposite $(-,-,0,0)$, and site $(+,-,0,0)$ has opposite $(-,+,0,0)$. The twenty-four sites pair into \textbf{twelve opposite-pairs}, and each such pair is called a \textbf{line}.

Each axis-pair block of four sites splits into two lines, the pair $(+,+)$ with $(-,-)$ and the pair $(+,-)$ with $(-,+)$. Six blocks times two lines gives $6 \times 2 = 12$ lines. A line carries the two tones moving \textbf{head-on} through the dock, one outbound along each end. The line, not the single site, is the unit the collision acts on.

\begin{definition}[Line]
A \emph{line} of a dock is a site together with its opposite, the pair of slots facing each other across the dock. The twenty-four sites form exactly twelve lines, and the knit's collision rewrites the tones one line at a time.
\end{definition}

\subsection{What one dock stores}

A dock holds \textbf{twenty-four ternary slots}, $s_0$ through $s_{23}$, one per site. Each slot carries a single \emph{tone}, a value in $\{-1, 0, +1\}$. The tone is the only substance in the theory. It is the vibe experience, read as pain, peace, and pleasure.

\begin{table}[ht]
\centering
\begin{tabular}{@{}llll@{}}
\toprule
value & name & feeling & meaning \\
\midrule
$-1$ & pain & aversion & a unit charge along that site, negative sign \\
$0$ & peace & neutral & empty, no charge in that slot \\
$+1$ & pleasure & attraction & a unit charge along that site, positive sign \\
\bottomrule
\end{tabular}
\caption{The single ternary tone and its three readings.}
\label{tab:the-tone}
\end{table}

Slot $s_k$ means the charge currently heading \textbf{out} through site $k$. That is the dock's entire state. Twenty-four ternary values, nothing else. No labels, no flags, no memory, no hidden variables. A dock is exactly its twenty-four tones.

The value order $-1 < 0 < +1$ is the \emph{lean}, the atom that tilts the tones. The lean is what makes peace the middle, the empty state from which a balanced $(+1, -1)$ pair can be born, and it is the seed from which charge and the direction of time later emerge. The dock's state space is therefore $3^{24}$, though in practice almost all of it is near-empty. A single moving charge is one slot equal to $+1$ with the other twenty-three at $0$.

\begin{principle}[No hidden state]
A dock is exactly its twenty-four ternary tones. It stores no edge weights, no stored connections, no internal memory. The single flat array of every dock's tones is the complete state of the universe.
\end{principle}

\subsection{The exact state of the whole world}

The whole world is \textbf{one flat array} of small integers, and there is nothing else. Every dock owns a run of twenty-four consecutive entries, one per site, laid end to end with the next dock's run.

\begin{table}[ht]
\centering
\begin{tabular}{@{}ll@{}}
\toprule
fact & value \\
\midrule
length & dock count times twenty-four \\
each entry & the tone of one dock in one site \\
allowed values & $-1$, $0$, $+1$ only \\
anything else stored & nothing, the array is the universe \\
\bottomrule
\end{tabular}
\caption{The entire state of the world as one array.}
\label{tab:the-world-array}
\end{table}

So the world is just the dock count times twenty-four trits. The tone of dock $c$ in site $d$ is the entry at position $24c + d$. To measure at a larger scale, grow the dock count. Never change a seed and never shrink a parameter for speed, since the initial state is a fixed deterministic function of the slot index, never a random draw.

\subsection{What a dock perceives, fill in and will out}

A dock does not store its connections. There are no edge weights and no stored links. The relation between two docks is just the shared face and the matching slots.

The \textbf{will} is the outgoing tone, slot $s_k$, the charge heading out through site $k$. The \textbf{fill} is the incoming tone, whatever the neighbours stream toward this dock. The \emph{site} itself is the perception across a direction, the one connection, and it is computed from the shared face, never stored.

Site $k$ of dock $A$ is the same shared face as one site of the neighbour reached along $k$. The two docks meet exactly there. $A$'s outgoing slot $s_k$ and the neighbour's slot for the same line of travel are the two ends of that one shared face. The perception is the face, the will is what $A$ pushes out through it, the fill is what arrives back.

\subsection{The charge, the conserved sum}

The familiar single tone of a dock, its net pain, peace, or pleasure, is just the \textbf{sum} of its twenty-four slots. The slots add which direction each bit of charge is heading. Sum them and the directions wash out, leaving the net local charge.

\begin{table}[ht]
\centering
\begin{tabular}{@{}ll@{}}
\toprule
object & definition \\
\midrule
dock & a here-now with twenty-four directed slots, one tone each \\
slot & the tone in one site at one dock \\
line & a site and its opposite, the pair of slots facing each other \\
dock tone & the sum of a dock's slots, its net local charge \\
charge & the sum of every slot over the whole mesh \\
\bottomrule
\end{tabular}
\caption{From a single slot up to the total charge.}
\label{tab:charge-ladder}
\end{table}

\begin{proposition}[Charge is conserved]
The collision of the knit preserves the sum of the tones on each line, and streaming only moves tones between docks without changing any value. Hence the total charge summed over every slot of every dock never changes. That conservation is the seed of identity.
\end{proposition}

\subsection{No rest slot is needed}

A dock has twenty-four sites and needs no rest slot. There is no zero-speed twenty-fifth site, and no home. A persistent self binds as a cluster of moving tones, held together by the emergent pressure of radiation in the active openness around it, not by any stationary slot. Its identity is the topological winding read at the resolution of the twenty-four sites. Neither the binding nor the identity needs a still center.

\begin{result}[Twenty-four, exactly]
The dock count is twenty-four sites, with no rest slot. The numbers $26$ and $80$ that appear elsewhere in the theory are the self's emergent direction-field counts, $3^3 - 1$ and $3^4 - 1$, a separate emergent structure. They are not dock tones. A dock has twenty-four sites.
\end{result}

\subsection{Matter and radiation, two motions of one tone}

There is one tone, not two. Matter and light are two \textbf{motions} of it. This is mass-energy equivalence built into the substance itself, since both are the same ternary stuff in motion.

\begin{table}[ht]
\centering
\begin{tabular}{@{}lll@{}}
\toprule
motion & what it is & physics name \\
\midrule
a charge bound and localized & a trapped, persistent lump of tone & matter, mass \\
a net-zero long-wavelength ripple & a propagating collective disturbance & light, sound, a phonon \\
\bottomrule
\end{tabular}
\caption{Two motions of the one tone.}
\label{tab:matter-radiation}
\end{table}

A bound tone is matter. The same tone waving is light. There are no separate fields.

\subsection{One beat, walked through}

A \emph{beat} is the discrete time step, and one beat is two steps, collide then stream. This two-step law is the \textbf{knit}.

\textbf{Step one, collide.} The dock reads its twenty-four slots and rewrites them in place by a fixed per-line table. The collision runs on each of the twelve lines independently. It conserves the per-line sum and is fully reversible. The move-types are a small fixed set. \emph{Inert}, where like charges coexist untouched. \emph{Hop}, where a lone charge crosses to the opposite slot. \emph{Create}, where peace brings out a balanced $(+1, -1)$ pair. \emph{Flip}, where such a pair reverses. \emph{Annihilate}, where a reversed pair returns to peace.

\textbf{Step two, stream.} Every tone moves to the neighbour its site points at. Slot $s_k$ of $A$ leaves along site $k$, arrives at the neighbour reached along $k$, and lands in that neighbour's slot for the \emph{same line of travel}, so it keeps going the same way. At the same instant $A$ receives, into each of its own slots, whatever its neighbours streamed toward it. Nothing decided direction. A charge heading one way simply kept going. This is momentum, ballistic at one dock per beat.

\begin{definition}[The knit]
The \emph{knit} is the state-law, collide then stream. Each beat the collision rewrites the tones in each dock line by line, by a fixed reversible charge-conserving table, then each tone streams one dock along its site, keeping its line of travel. The knit interlaces and advances beat by beat, weaving the fabric of the mesh.
\end{definition}

\subsection{A worked example}

Start with dock $A$ holding $s_3 = +1$ and all other twenty-three slots at $0$, a single $+1$ charge heading out along site $3$. All neighbours are empty.

\begin{enumerate}
\item \textbf{Collide at $A$.} The line $(3, \mathrm{opposite}(3))$ reads $(+1, 0)$. The per-line table hops the lone charge to $(0, +1)$. The charge has crossed the dock.
\item \textbf{Stream.} The tone leaves $A$ along its site and arrives at the neighbour, landing in that neighbour's straight-ahead slot. $A$ empties, the neighbour now holds one incoming $+1$.
\item \textbf{Repeat.} Beat after beat the $+1$ streams in a straight line, one dock per beat, never spreading. That is a \textbf{massless} particle.
\item \textbf{Add mass.} Turn on a mixing collision, a scatter route, and the charge begins turning into other slots. Its track zig-zags and its effective speed drops. That is \textbf{mass}.
\item \textbf{Pair creation.} Start instead from an all-peace region with the create move on. A $+1$ and a $-1$ are emitted into the two slots of one line, net zero, and they fly apart. That is \textbf{pair production}. Throughout, the sum over all docks' slots never changes. Charge is conserved.
\end{enumerate}

\subsection{What this fixes in intuition}

Nothing moves through space. Space is the dock-graph. The only motion possible is a tone-pattern re-drawn one dock over each beat, a glider.

\begin{itemize}
\item \textbf{Momentum} is which slot a charge lives in.
\item \textbf{Mass} is how much the collision mixes slots.
\item \textbf{Pair creation} is the create move of the collision.
\item \textbf{Conservation} is simply that the collision never changes the sum.
\end{itemize}

One dock, twenty-four ternary slots, two steps a beat. That is the whole machine. The mesh grows at its \emph{wake}, new docks born at peace along the moving frontier, and the direction of time, the \emph{arrow}, emerges from that growth acting through the lean. Everything in later sections is what this machine does at scale.
